\section{The instrument (Layer 1 --- measurement validity)}
\label{sec:instrument}

\aidraft{%
The estimator is a two-state weighted statistical jump model \citep{bemporad2018,nystrup2020eswa}
fit to three return-only features --- exponentially weighted downside deviation (10-day half-life)
and Sortino ratios at 20- and 60-day half-lives --- with a jump penalty $\lambda$ that enforces
state persistence.
We evaluate it first purely as a measurement device, before any economic use, on three properties
a deployable instrument must have: stability under reasonable estimation choices, agreement with an
external ground truth, and live computability.%
}

\paragraph{Stability.}
\aidraft{%
The property that motivated the whole program is reproducibility of the label under
training-window perturbation.
Shifting the training start by $\pm 2$ years --- a first-order analyst choice that leaves competing
estimators disagreeing about the state on a fifth of all days --- changes the jump-model label on
$0.0\%$ of out-of-sample days: \stabJM{}/\stabJM{} agreement at both the $-2$y and $+2$y shifts.
An HMM ensemble on the identical data and features achieves only \stabHMM{}.
The label switches \switchesYr{} times per year and partitions the 1990--2026 out-of-sample period
into \nEpisodes{} episodes, each a genuine high-volatility environment (calm-state annualized
volatility ${\approx}\,\calmVol$, stressed-state ${\approx}\,\stressVol$); $\lambda$ selects to the
interior of its search grid rather than an edge.
Estimation instability --- the practical reason HMM regime labels are hard to trust in production
--- is, for this estimator on this data, solved.%
}
\reviewnote{Adam: 1.000 is under-argued as it stands. A constant label is also perfectly stable ---
so pair stability with skill; harden the perturbation ($\pm$5y, bootstrap-resample, param jitter)
and show the degradation curve; add a deterministic vol-threshold baseline next to the stochastic
HMM. This is the referee's first objection. (See NOTES.md instrument to-do.)}

\paragraph{External agreement, and an honest boundary.}
\aidraft{%
Validated against ex-post bear-market datings, the label catches \bearsCaught{} declines of
$-15\%$ or more, at a median recognition lag of \detLag{} trading days from the ex-post peak.
It is, precisely, a \emph{volatility-state} detector rather than a \emph{bear-market} detector: it
fires on high-volatility environments whether or not a sustained decline follows (hence a moderate
precision against directional bear datings), and it misses fast crashes that resolve before its
persistence-seeking filter can accumulate evidence --- 1998 (LTCM) and 2018-Q4 are sub-threshold at
the deployed $\lambda$.
We state this as the instrument's boundary of use rather than a defect to be tuned away: the
\detLag{}-day lag and the fast-crash misses are exactly the constraints the mechanism section shows
are binding for any causal detector, and they define what the monitor (Section~\ref{sec:monitor})
can and cannot claim.%
}

\paragraph{Live computability.}
\aidraft{%
The instrument runs to the present.
Ken French publishes with a one-to-two-month lag, so the live tail is spliced from an SPY proxy;
the splice passes its gate at \spliceCorr{} return correlation, and the composite label agrees with
the frozen research label on \overlapAgree{} of the \overlapDays{}-day overlap.
There is no vintage-lookahead: the filter's state at time $t$ uses only data through $t$, and the
live label at $t$ is the same object the frozen backtest would have produced.%
}

\aidraft{%
These three properties make the label a measurement instrument worth building on.
The remaining sections ask what happens when a decision is conditioned on it --- and find,
uniformly, that the measurement's quality does not transfer to decision value.%
}
